\documentclass[conference]{IEEEtran}
\IEEEoverridecommandlockouts

\usepackage{cite}
\usepackage{amsmath,amssymb,amsfonts}
\usepackage{algorithmic}
\usepackage{graphicx}
\usepackage{textcomp}
\usepackage{xcolor}
\usepackage{booktabs}
\usepackage{multirow}
\usepackage{url}
\usepackage{microtype}

\def\BibTeX{{\rm B\kern-.05em{\sc i\kern-.025em b}\kern-.08em
    T\kern-.1667em\lower.7ex\hbox{E}\kern-.125emX}}

\begin{document}

\title{Psychological-State-Aware Conversational AI: A Real-Time Psychologically-Adaptive Conversational AI Framework\\
with Multimodal Emotion Fusion}

\author{%
  \IEEEauthorblockN{Sravana Kumar K} 
  \IEEEauthorblockA{\textit{Dept.\ of CSE (AI\&ML)}\\
    Vishnu Institute of Technology\\
    Bhimavaram, India\\
    sravanakumar.k@vishnu.edu.in} 
  \and
  \IEEEauthorblockN{Sai Pavan Kumar Devisetti} 
  \IEEEauthorblockA{\textit{Dept.\ of CSE (AI\&ML)}\\
    Vishnu Institute of Technology\\
    Bhimavaram, India\\
    22pa1a4231@vishnu.edu.in} 
  \and
  \IEEEauthorblockN{Giridha Sri Meghana.V}
  \IEEEauthorblockA{\textit{Dept.\ of CSE (AI\&ML)}\\
    Vishnu Institute of Technology\\
    Bhimavaram, India\\
    22pa1a4237@vishnu.edu.in} 
  \and
  \IEEEauthorblockN{V Sujan Ajitesh Bayya} 
  \IEEEauthorblockA{\textit{Dept.\ of CSE (AI\&ML)}\\
    Vishnu Institute of Technology\\
    Bhimavaram, India\\
    22pa1a4212@vishnu.edu.in} 
  \and
  \IEEEauthorblockN{Shanmukha Suryanarayan Chinta} 
  \IEEEauthorblockA{\textit{Dept.\ of CSE (AI\&ML)}\\
    Vishnu Institute of Technology\\
    Bhimavaram, India\\
    22pa1a4224@vishnu.edu.in} 
}

\maketitle

\begin{abstract}
Conversational AI systems are increasingly expected to operate in emotionally complex and psychologically sensitive contexts. Yet most deployed dialogue agents remain agnostic to the psychological state of the user, relying solely on linguistic content and ignoring the rich affective signals present in speech prosody and acoustic events. In this paper, we present a Psychological-State-Aware Conversational AI (PSACA) system that integrates multimodal emotion recognition, adaptive psychological trend modeling, and event-driven long-term memory retrieval to enable empathic, contextually coherent dialogue. The system constructs a four-dimensional psychological state vector comprising valence, arousal, stress, and cognitive clarity, derived from the fusion of text-based emotion analysis, speech-based emotion recognition, and acoustic event detection. Psychological trends are modeled via a layered pipeline of instantaneous signals, Exponential Moving Average (EMA) smoothing, regression-based trajectory estimation, and hysteresis gating to prevent spurious state oscillation. Long-term memory is maintained through an event-driven vector database, while a psychologically informed prompting strategy conditions a large language model (LLM) on trend trajectories rather than instantaneous states. Across a controlled evaluation of 100 conversational turns, our system achieves Trend Stability Scores (TSS) of 0.93, 0.91, 0.89, and 0.95 for valence, arousal, stress, and clarity respectively, a Memory Storage Rate of 34\%, and an Emotional Shift Detection Rate of 17\%. These results demonstrate the viability of continuous psychological state tracking in real-time conversational AI and highlight directions for future empathic dialogue systems.
\end{abstract}

\begin{IEEEkeywords}
affective computing, conversational AI, multimodal emotion recognition, psychological state modeling, adaptive trend modeling, empathic dialogue, memory-augmented dialogue systems, large language models
\end{IEEEkeywords}

% ============================================================
\section{Introduction}
% ============================================================

Human communication is inherently affective. When people engage in conversation, they simultaneously convey information, emotional state, cognitive load, and psychological disposition through multiple channels---words, prosody, rhythm, and acoustic texture \cite{picard2000affective, schuller2018age}. While natural language understanding has advanced substantially through transformer-based architectures and large language models (LLMs) \cite{brown2020language, openai2023gpt4}, the dominant paradigm in conversational AI remains fundamentally reactive: the system responds to what is said, not to how the speaker feels or how their psychological state is evolving over time.

This limitation becomes particularly consequential in high-stakes conversational domains such as mental health support, coaching, tutoring, and crisis intervention, where misreading or ignoring a user's emotional trajectory can undermine trust, escalate distress, or produce contextually inappropriate responses \cite{morris2018towards, fitzpatrick2017delivering}. A conversational agent that cannot distinguish a user who is calmly inquiring about career options from one who is anxious and seeking reassurance will fail to adapt its tone, pacing, and content accordingly.

The challenge of building psychologically aware conversational AI involves several interrelated sub-problems. First, emotional and psychological signals must be recognized from multiple modalities simultaneously, since any single channel---text, voice, or audio events---provides only a partial and noisy view of the user's internal state \cite{poria2017review, sebe2005multimodal}. Second, the resulting state estimates must be smoothed and trended over time to distinguish transient fluctuations from genuine psychological shifts \cite{cowie2001emotion}. Third, relevant context from prior interactions must be selectively retrieved to inform the current response without overwhelming the model's context window \cite{zhong2022dialoglm}. Finally, the LLM generating the response must be conditioned on this rich psychological context in a structured and principled manner.

In this paper, we present \textit{PSACA} (Psychological-State-Aware Conversational AI), a modular end-to-end system that addresses each of these sub-problems through a cohesive architecture. Our core contributions are as follows:

\begin{itemize}
    \item A \textbf{multimodal psychological fusion engine} that combines text emotion classification, speech-based emotion recognition, and acoustic event detection into a unified four-dimensional psychological state vector: valence ($v$), arousal ($a$), stress ($s$), and cognitive clarity ($c$).
    \item An \textbf{adaptive trend modeling pipeline} employing instantaneous signals, Exponential Moving Average (EMA) smoothing, linear regression-based trajectory estimation, and hysteresis gating for robust, oscillation-resistant state tracking.
    \item An \textbf{event-driven long-term memory subsystem} backed by vector storage that selectively archives and retrieves semantically relevant conversational episodes based on psychological salience.
    \item A \textbf{trend-conditioned LLM prompting strategy} that informs the generative model with psychological trajectory data rather than raw instantaneous state estimates.
    \item \textbf{Interaction-mode-adaptive text-to-speech (TTS)} synthesis that modulates prosodic parameters in response to the inferred interaction mode (de-escalation, clarification, flow, or neutral).
\end{itemize}

We evaluate PSACA over a longitudinal sequence of 100 conversational turns spanning diverse psychological topics and demonstrate quantitatively that the system maintains stable, coherent psychological state estimates while accurately detecting meaningful emotional transitions.

The remainder of this paper is organized as follows. Section~\ref{sec:related} reviews related work. Section~\ref{sec:methodology} details the system methodology. Section~\ref{sec:architecture} describes the system architecture. Section~\ref{sec:fusion} presents the multimodal fusion design. Section~\ref{sec:trends} details the trend modeling approach. Section~\ref{sec:memory} covers the memory architecture. Section~\ref{sec:evaluation} describes the evaluation framework. Section~\ref{sec:results} presents results. Section~\ref{sec:discussion} discusses implications. Section~\ref{sec:limitations} discusses limitations, and Section~\ref{sec:conclusion} concludes.

% ============================================================
\section{Related Work}
\label{sec:related}
% ============================================================

\subsection{Affective Computing and Emotion Recognition}

The field of affective computing, introduced by Picard \cite{picard2000affective}, established the foundational premise that machines should recognize, interpret, and simulate human affect. Early work focused primarily on unimodal signals, including facial expressions \cite{ekman1978facial}, speech prosody \cite{schuller2003hidden}, and physiological signals \cite{healey2005detecting}. More recent work has shifted toward multimodal fusion, combining linguistic, acoustic, and visual cues to achieve more robust emotion recognition \cite{poria2017review, zadeh2018multimodal, tsai2019multimodal}.

Text-based emotion classification has benefited substantially from pre-trained language models. Models such as BERT \cite{devlin2019bert} and its affective variants have shown strong performance on sentiment and emotion classification tasks \cite{sun2019utilizing}. Speech-based emotion recognition has similarly advanced through deep neural architectures applied to mel-spectrograms and acoustic feature sets such as eGeMAPS \cite{eyben2016geneva, schuller2018age}.

\subsection{Psychological State Modeling}

The dimensional model of emotion, comprising valence and arousal axes, provides a compact and theoretically grounded framework for representing continuous emotional states \cite{russell1980circumplex}. Extensions to this model have incorporated additional dimensions such as dominance \cite{mehrabian1996pleasure} and cognitive load \cite{sweller1988cognitive}. Our system adopts a four-dimensional representation---valence, arousal, stress, and clarity---that captures both affective and cognitive aspects of psychological state relevant to empathic conversation.

Temporal modeling of emotional state has received comparatively less attention. Cowie et al. \cite{cowie2001emotion} argued that emotion recognition must account for temporal dynamics rather than treating each utterance as an independent classification problem. EMA-based smoothing has been applied in speech processing \cite{rabiner1993fundamentals} and recently explored for affective state tracking \cite{metallinou2013tracking}.

\subsection{Memory-Augmented Dialogue Systems}

Long-term memory in dialogue systems has been approached through various mechanisms including explicit memory slots \cite{sukhbaatar2015end}, retrieval-augmented generation (RAG) \cite{lewis2020retrieval}, and dense vector retrieval \cite{karpukhin2020dense}. Recent work on persona-consistent and memory-grounded dialogue \cite{xu2021beyond, zhong2022dialoglm} has demonstrated the value of selectively retrieving and grounding on prior conversational episodes. Our approach extends this line of work by conditioning memory retrieval on psychological salience in addition to semantic relevance.

\subsection{Empathic Conversational Agents}

Woebot \cite{fitzpatrick2017delivering} and similar systems demonstrated that rule-based empathic conversational agents can provide meaningful psychological support. More recent systems have employed LLMs for open-domain empathic dialogue \cite{rashkin2019towards, majumder2020mime}. However, these systems typically lack continuous psychological state tracking and rely on turn-level rather than trajectory-level affective modeling. PSACA addresses this gap by maintaining a continuously updated psychological state estimate that informs every stage of response generation.

\subsection{Adaptive TTS in Dialogue}

Expressive and adaptive TTS systems have explored prosody control \cite{shen2018natural, ren2020fastspeech} and emotion-conditioned speech synthesis \cite{lei2022msemotts}. The integration of dialogue-state-aware TTS with affective conversational agents remains an underexplored area, which our system addresses through interaction-mode-conditioned speech synthesis.

% ============================================================
\section{Methodology}
\label{sec:methodology}
% ============================================================

\subsection{Problem Formulation}

Let $\mathcal{C} = \{(u_1, r_1), (u_2, r_2), \ldots, (u_T, r_T)\}$ denote a conversational session of $T$ turns, where $u_t$ is the user utterance at turn $t$ and $r_t$ is the system response. Each utterance $u_t$ is associated with a multimodal observation tuple $\mathbf{o}_t = (\mathbf{x}^{\text{txt}}_t, \mathbf{x}^{\text{aud}}_t, \mathbf{x}^{\text{evt}}_t)$, where $\mathbf{x}^{\text{txt}}_t$ denotes text features, $\mathbf{x}^{\text{aud}}_t$ denotes acoustic/speech features, and $\mathbf{x}^{\text{evt}}_t$ denotes detected audio events.

The goal of PSACA is to maintain, at each turn $t$, a psychological state estimate $\boldsymbol{\psi}_t = (v_t, a_t, s_t, c_t)$ where $v_t \in [-1, 1]$ is valence, $a_t \in [0, 1]$ is arousal, $s_t \in [0, 1]$ is stress, and $c_t \in [0, 1]$ is cognitive clarity. The system then conditions response generation on the trend trajectory $\{\boldsymbol{\psi}_{t-k}, \ldots, \boldsymbol{\psi}_t\}$ and relevant retrieved memory episodes, rather than solely on $\boldsymbol{\psi}_t$.

\subsection{Multimodal Feature Extraction}

Text features are extracted by passing $u_t$ through a transformer-based emotion classifier that outputs a probability distribution over emotion categories (e.g., joy, sadness, anger, fear, surprise, disgust, neutral), from which dimensional state estimates are derived via affective mappings. Acoustic features are extracted from the audio waveform of $u_t$ using a pre-trained speech emotion recognition model operating on mel-spectrogram representations, yielding an independent emotion probability distribution. Acoustic events (e.g., laughter, crying, silence, hesitation) are detected using a lightweight audio event classifier operating on short-time Fourier transform features.

\subsection{Psychological State Fusion}

The multimodal fusion module combines the three modality-level estimates into a unified psychological state vector. Denoting the text-derived dimensional estimates as $\boldsymbol{\psi}^{\text{txt}}_t$, the audio-derived estimates as $\boldsymbol{\psi}^{\text{aud}}_t$, and the event-derived modifiers as $\Delta\boldsymbol{\psi}^{\text{evt}}_t$, the fused state is computed as:

\begin{equation}
\boldsymbol{\psi}^{\text{fused}}_t = \alpha \cdot \boldsymbol{\psi}^{\text{txt}}_t + \beta \cdot \boldsymbol{\psi}^{\text{aud}}_t + \gamma \cdot \Delta\boldsymbol{\psi}^{\text{evt}}_t
\label{eq:fusion}
\end{equation}

where $\alpha, \beta, \gamma \in [0,1]$ are modality confidence weights satisfying $\alpha + \beta + \gamma = 1$, dynamically adjusted based on signal reliability scores. Valence is mapped to $[-1, 1]$ while arousal, stress, and clarity are mapped to $[0, 1]$ via sigmoid normalization.

\subsection{Adaptive Trend Modeling}

A core methodological contribution of PSACA is its layered approach to psychological trend estimation. Rather than acting on instantaneous state estimates, which are inherently noisy, the system maintains four parallel representations of each psychological dimension.

\textbf{Instantaneous signal} $\psi^{\text{inst}}_t$ is the direct output of the fusion module for turn $t$.

\textbf{Exponential Moving Average (EMA)} provides first-order smoothing:
\begin{equation}
\psi^{\text{ema}}_t = \lambda \cdot \psi^{\text{inst}}_t + (1 - \lambda) \cdot \psi^{\text{ema}}_{t-1}
\label{eq:ema}
\end{equation}
where $\lambda \in (0, 1)$ is the smoothing factor (empirically set to $\lambda = 0.3$ in our system).

\textbf{Regression-based trend} $\psi^{\text{reg}}_t$ is estimated by fitting a linear regression model over a sliding window of the $W$ most recent EMA values, providing a trajectory estimate:
\begin{equation}
\psi^{\text{reg}}_t = \hat{\beta}_0 + \hat{\beta}_1 \cdot t
\label{eq:regression}
\end{equation}
where $\hat{\beta}_0$ and $\hat{\beta}_1$ are ordinary least-squares coefficients estimated over the window $[t-W+1, t]$.

\textbf{Hysteresis gating} prevents rapid oscillation by requiring a state transition to exceed a threshold $\delta$ relative to the last committed state before updating the trend label:
\begin{equation}
\psi^{\text{gate}}_t = \begin{cases}
\psi^{\text{reg}}_t & \text{if } |\psi^{\text{reg}}_t - \psi^{\text{gate}}_{t-1}| > \delta \\
\psi^{\text{gate}}_{t-1} & \text{otherwise}
\end{cases}
\label{eq:hysteresis}
\end{equation}

\subsection{Trend Stability Score}

To quantify the smoothness and coherence of the tracked psychological state, we define the Trend Stability Score (TSS) for a dimension as:

\begin{equation}
\text{TSS} = 1 - \frac{1}{T} \sum_{t=1}^{T} |\psi^{\text{inst}}_t - \psi^{\text{ema}}_t|
\label{eq:tss}
\end{equation}

A TSS approaching 1.0 indicates that the instantaneous signals closely track the smoothed EMA, reflecting low noise and high signal coherence. TSS values substantially below 1.0 indicate high volatility in the instantaneous signal relative to the smoothed trend.

\subsection{Memory Storage Rate}

At each turn, the memory manager evaluates the psychological salience of the current episode and decides whether to archive it to long-term memory. The Memory Storage Rate is defined as:

\begin{equation}
\text{MSR} = \frac{N_{\text{stored}}}{N_{\text{total}}}
\label{eq:msr}
\end{equation}

where $N_{\text{stored}}$ is the number of turns archived to memory and $N_{\text{total}}$ is the total number of turns evaluated. A selective MSR reflects a policy that prioritizes salient or novel episodes over redundant routine exchanges.

\subsection{Trend-Conditioned LLM Prompting}

At response generation time, the LLM client constructs a structured prompt incorporating: (i) the bounded rolling conversational context from the current session, (ii) the four-dimensional psychological trend trajectory over the most recent $K$ turns, (iii) the inferred interaction mode, and (iv) retrieved memory episodes. This trend-conditioned prompt enables the LLM to generate responses that are sensitive not only to what was said, but to how the user's psychological state has been evolving, supporting more empathic and contextually appropriate dialogue.

% ============================================================
\section{System Architecture}
\label{sec:architecture}
% ============================================================

The PSACA system is composed of nine tightly integrated modules organized into a pipeline that processes each user utterance end-to-end, from audio input to spoken response. Figure~\ref{fig:architecture} illustrates the complete system architecture.

\begin{figure}[htbp]
\centering
\includegraphics[width=\linewidth]{System Architecture.png}
\caption{PSACA System Architecture. The pipeline processes raw audio through ASR, multimodal emotion analysis, psychological fusion, trend modeling, memory retrieval, LLM prompting, and adaptive TTS synthesis.}
\label{fig:architecture}
\end{figure}

\textbf{realtime\_conversational\_ai} serves as the system orchestrator, coordinating the full pipeline across turns and managing the session lifecycle.

\textbf{services} provides three co-located services: Automatic Speech Recognition (ASR) for transcribing user audio, text emotion analysis for classifying the linguistic content of utterances, and audio emotion analysis for classifying speech-prosody-level emotion from the acoustic waveform.

\textbf{fusion} implements the multimodal psychological fusion engine described in Section~\ref{sec:fusion}, combining the three modality-level emotion estimates into the unified psychological state vector $\boldsymbol{\psi}_t$.

\textbf{emotional\_trends} maintains the adaptive trend modeling pipeline described in Section~\ref{sec:methodology}, computing EMA, regression-based trend, and hysteresis-gated state for each psychological dimension across turns.

\textbf{memory\_store} provides vector database storage for long-term conversational episodes, supporting dense nearest-neighbor retrieval by semantic embedding similarity.

\textbf{memory\_orchestrator} implements event-driven memory retrieval logic, triggering retrieval when the current turn's psychological profile and topic embedding exceed similarity thresholds relative to archived episodes.

\textbf{memory\_manager} performs topic detection using embedding clustering and applies a configurable memory policy that scores each turn for psychological salience and novelty before committing it to storage.

\textbf{llm\_client} constructs trend-conditioned prompts and interfaces with GPT-4o-mini for response generation, injecting the psychological trajectory, interaction mode, and retrieved memories as structured context.

\textbf{tts\_azure} implements interaction-mode-adaptive TTS synthesis, modulating speaking rate, pitch variation, and pause distribution based on the current inferred interaction mode (de-escalation, clarification, flow, or neutral).

\textbf{conversation\_state} maintains a bounded rolling context window over the current session, ensuring the prompt size remains within LLM context limits while preserving the most recent and most psychologically salient exchanges.

% ============================================================
\section{Multimodal Psychological Fusion}
\label{sec:fusion}
% ============================================================

The fusion module is responsible for combining evidence from three distinct signal streams into a coherent psychological state estimate. Figure~\ref{fig:fusion} illustrates the fusion and emotional trend modeling pipeline.

\begin{figure}[htbp]
\centering
\includegraphics[width=\linewidth]{Fusion and emotion trend.png}
\caption{Multimodal Fusion and Emotional Trend Modeling. Text emotion, audio emotion, and acoustic event signals are fused into a four-dimensional psychological state vector, which is subsequently processed through EMA smoothing, regression-based trend estimation, and hysteresis gating.}
\label{fig:fusion}
\end{figure}

\subsection{Modality-Level Processing}

Text emotion analysis is performed by a transformer encoder fine-tuned on an affective dialogue corpus, producing logit scores over seven emotion categories. These categorical probabilities are mapped to the valence-arousal space using an affective mapping matrix derived from established dimensional-categorical correspondences \cite{russell1980circumplex}. Stress is estimated as a weighted combination of high-arousal negative emotion probabilities (anger, fear, disgust), while clarity is derived inversely from the entropy of the emotion distribution---high distributional uncertainty signals cognitive confusion or mixed affect.

Speech emotion recognition operates directly on the acoustic waveform, extracting mel-frequency cepstral coefficients (MFCCs), fundamental frequency contours, and energy statistics. A lightweight convolutional recurrent neural network produces a per-utterance emotion probability vector that is similarly mapped to dimensional estimates. The acoustic emotion estimate provides a complementary view of psychological state, particularly valuable when linguistic content is affectively neutral or ambiguous.

Acoustic event detection identifies salient non-linguistic events (laughter, crying, sighing, prolonged silence, speech disfluencies) using a multi-label classifier. Detected events apply structured modifications to the fused state vector: for example, a detected sigh increases the stress estimate and decreases arousal, while laughter shifts valence positively and raises arousal.

\subsection{Confidence-Weighted Fusion}

Modality confidence weights $\alpha$, $\beta$, and $\gamma$ in Equation~\eqref{eq:fusion} are dynamically estimated per turn. The text modality confidence is proportional to the speech recognition confidence score and the entropy of the text emotion distribution. The audio modality confidence is proportional to the signal-to-noise ratio of the acoustic stream. The event modality weight $\gamma$ is set to a small base value and increased when high-confidence events are detected.

\subsection{Interaction Mode Classification}

The fused state vector $\boldsymbol{\psi}^{\text{fused}}_t$ is used to classify the current interaction into one of four modes:

\begin{itemize}
    \item \textbf{De-escalation}: high stress ($s_t > \theta_s$) or strongly negative valence ($v_t < \theta_v^-$), indicating the system should prioritize calming and grounding responses.
    \item \textbf{Clarification}: low clarity ($c_t < \theta_c$), indicating cognitive confusion and prompting the system to provide structured, clear information.
    \item \textbf{Flow}: high clarity and moderate positive valence, indicating productive and engaged dialogue.
    \item \textbf{Neutral}: default mode when no specialized interaction adaptation is warranted.
\end{itemize}

Thresholds $\theta_s$, $\theta_v^-$, and $\theta_c$ are empirically calibrated on held-out conversational data.

% ============================================================
\section{Trend Modeling}
\label{sec:trends}
% ============================================================

A fundamental challenge in psychological state tracking is distinguishing meaningful state transitions from transient noise. A user who briefly mentions a stressful event does not necessarily exhibit a sustained increase in stress; conversely, a gradual drift in valence across many turns may signal a genuine psychological shift. The PSACA trend modeling pipeline addresses this challenge through its layered architecture.

\subsection{EMA Smoothing}

EMA smoothing, defined in Equation~\eqref{eq:ema}, provides a computationally efficient first-order low-pass filter over the instantaneous state estimates. With $\lambda = 0.3$, recent observations receive 30\% weight while the running history receives 70\% weight, resulting in a smoothed estimate that is responsive to genuine change while attenuating transient spikes.

\subsection{Regression-Based Trajectory}

The regression-based trend estimate provides a forward-looking trajectory rather than a simple running average. By fitting a linear model over the $W = 10$ most recent EMA values, the trend module captures directional momentum: whether the psychological state is improving, deteriorating, or stable. This trajectory information is conveyed to the LLM client as a structured description (e.g., ``valence is trending upward, stress is declining'') that guides response generation toward trajectory-appropriate behaviors.

\subsection{Hysteresis Gating}

Without hysteresis, minor fluctuations in the regression trend could trigger frequent mode switches, producing inconsistent system behavior. The hysteresis gate, defined in Equation~\eqref{eq:hysteresis}, requires that a dimension's trend deviate from its last committed value by at least $\delta = 0.1$ (on the normalized scale) before the trend label is updated. This mechanism substantially reduces interaction mode churn while preserving sensitivity to genuine psychological transitions.

% ============================================================
\section{Memory Architecture}
\label{sec:memory}
% ============================================================

\subsection{Memory Store}

Long-term memory is implemented as a dense vector store (following a Qdrant-style design) where each stored episode is represented by a semantic embedding of the conversational turn concatenated with a psychological state fingerprint. Retrieval is performed via approximate nearest-neighbor search, combining semantic similarity with psychological state similarity through a weighted composite distance metric.

\subsection{Memory Orchestrator}

The memory orchestrator implements event-driven retrieval rather than querying the memory store at every turn. Retrieval is triggered when one or more of the following conditions are met: (i) the current topic embedding diverges substantially from the recent topic trajectory, indicating a topic shift; (ii) the inferred psychological state exhibits a sharp transition (as detected by the hysteresis gating module); or (iii) the current turn is classified as psychologically salient by the memory manager. This selective retrieval policy minimizes LLM prompt overhead while ensuring relevant historical context is surfaced when most needed.

\subsection{Memory Manager}

The memory manager performs two primary functions. First, it performs real-time topic detection using a lightweight sentence embedding model that clusters utterances into high-level topic categories. Second, it applies a configurable memory policy that scores each candidate episode for storage based on psychological salience (measured as deviation from the session mean psychological state), topic novelty, and conversational significance. Episodes that score above a calibrated threshold are archived to the vector store.

% ============================================================
\section{Evaluation Framework}
\label{sec:evaluation}
% ============================================================

\subsection{Evaluation Protocol}

We evaluated PSACA over a controlled sequence of 100 conversational turns spanning five psychological topic domains: career and startup challenges, interpersonal relationships, stress and anxiety, personal growth, and daily life. Each turn was processed by the full system pipeline. Evaluation focused on the quality and stability of psychological state tracking, the appropriateness of interaction mode classification, memory system behavior, and emotional shift detection.

\subsection{Metrics}

\textbf{Trend Stability Score (TSS)} (Equation~\eqref{eq:tss}) measures the coherence between instantaneous and EMA-smoothed estimates per psychological dimension.

\textbf{Interaction Mode Distribution} measures the proportion of turns assigned to each interaction mode, providing a behavioral profile of the system across the evaluation session.

\textbf{Memory Storage Rate (MSR)} (Equation~\eqref{eq:msr}) measures the selectivity of the memory policy.

\textbf{Emotional Shift Detection Rate (ESDR)} measures the proportion of turns at which the hysteresis gate detected a committed state transition exceeding threshold $\delta$ in at least one psychological dimension.

\textbf{Psychological State Variance Reduction (PSVR)} measures the proportional reduction in per-dimension variance between instantaneous estimates and EMA-smoothed estimates:
\begin{equation}
\text{PSVR} = 1 - \frac{\text{Var}(\psi^{\text{ema}}_{1:T})}{\text{Var}(\psi^{\text{inst}}_{1:T})}
\label{eq:psvr}
\end{equation}

\textbf{Stability Index Ratio (SIR)} measures the ratio of EMA-smoothed signal variance to raw signal variance, where lower values indicate greater smoothing effectiveness.

% ============================================================
\section{Results}
\label{sec:results}
% ============================================================

\subsection{Trend Stability Scores}

Table~\ref{tab:tss} reports TSS values for each psychological dimension across the 100-turn evaluation session. All dimensions achieve TSS values above 0.89, demonstrating strong alignment between instantaneous signals and EMA-smoothed estimates.

\begin{table}[htbp]
\centering
\caption{Trend Stability Scores (TSS) Across 100 Conversational Turns}
\label{tab:tss}
\begin{tabular}{lcc}
\toprule
\textbf{Dimension} & \textbf{TSS} & \textbf{Interpretation} \\
\midrule
Valence   & 0.93 & Very high stability \\
Arousal   & 0.91 & High stability \\
Stress    & 0.89 & High stability \\
Clarity   & 0.95 & Very high stability \\
\midrule
\textbf{Mean} & \textbf{0.92} & \\
\bottomrule
\end{tabular}
\end{table}

Figure~\ref{fig:valence} illustrates the valence trajectory, showing the close tracking between instantaneous, EMA, and regression-based estimates across turns.

\begin{figure}[htbp]
\centering
\includegraphics[width=\linewidth]{100turns_figure_1_valence.png}
\caption{Valence trajectory across 100 conversational turns. Instantaneous signal (blue), EMA-smoothed estimate (orange), and regression-based trend (green) are shown. The close tracking between instantaneous and EMA signals is reflected in the high TSS of 0.93.}
\label{fig:valence}
\end{figure}

\begin{figure}[htbp]
\centering
\includegraphics[width=\linewidth]{100turns_figure_4_tss.png}
\caption{Trend Stability Score (TSS) per psychological dimension across the 100-turn evaluation session. Clarity achieves the highest TSS (0.95) while Stress exhibits the most variability (TSS = 0.89), consistent with the inherently volatile nature of stress responses in conversational contexts.}
\label{fig:tss}
\end{figure}

\subsection{Interaction Mode Distribution}

Figure~\ref{fig:modes} shows the distribution of interaction modes across the evaluation session.

\begin{figure}[htbp]
\centering
\includegraphics[width=\linewidth]{100turns_figure_2_modes.png}
\caption{Interaction mode distribution across 100 conversational turns. Neutral mode predominates (44\%), followed by Flow (26\%), De-escalation (18\%), and Clarification (12\%). The distribution reflects a realistic balance of psychologically stable and psychologically demanding exchanges.}
\label{fig:modes}
\end{figure}

The mode distribution---Neutral (44\%), Flow (26\%), De-escalation (18\%), Clarification (12\%)---reflects a realistic representation of conversational variability. The relatively high incidence of de-escalation mode (18\%) is consistent with the psychological topic domains covered, which included stress, anxiety, and interpersonal challenges.

\subsection{Memory System Performance}

The memory manager archived 34\% of evaluated turns (Memory Storage Rate = 0.34), demonstrating selective rather than indiscriminate storage. Figure~\ref{fig:memory} illustrates the memory storage pattern across turns, showing clustering of archived episodes around topic transitions and emotionally salient moments.

\begin{figure}[htbp]
\centering
\includegraphics[width=\linewidth]{100turns_figure_3_memory.png}
\caption{Memory storage events across 100 conversational turns. Archived turns (marked) are concentrated around topic transitions and moments of elevated psychological salience, demonstrating the selectivity of the memory policy.}
\label{fig:memory}
\end{figure}

\subsection{Emotional Shift Detection}

The system detected committed emotional shifts in 17\% of turns (Emotional Shift Detection Rate = 0.17). Figure~\ref{fig:shifts} illustrates detected shifts across the evaluation session, showing their distribution relative to topic boundaries and interaction mode transitions.

\begin{figure}[htbp]
\centering
\includegraphics[width=\linewidth]{100turns_figure_6_shifts.png}
\caption{Emotional shift detection events across 100 conversational turns. Detected shifts (17\% of turns) are concentrated at topic boundaries and during transitions between interaction modes, confirming the temporal validity of the hysteresis-gated detection mechanism.}
\label{fig:shifts}
\end{figure}

\subsection{Psychological State Variance Reduction}

Applying Equation~\eqref{eq:psvr}, we find that EMA smoothing reduces per-dimension signal variance by 38.4\% on average (valence: 41.2\%, arousal: 37.8\%, stress: 34.6\%, clarity: 40.0\%). This variance reduction quantifies the noise-suppression effectiveness of the EMA layer and underscores the value of trend-based rather than instantaneous state conditioning for LLM prompting.

\subsection{Topic Distribution}

The memory manager's topic classifier assigned turns across five domains: Career/Startup (30\%), Relationships (21\%), Stress/Anxiety (18\%), Personal Growth (17\%), and Daily Life (14\%). This distribution is consistent with the conversational scenarios covered in the evaluation session and demonstrates the topic classifier's ability to discriminate across semantically diverse domains.

\subsection{Stability Index Comparison}

Table~\ref{tab:sir} reports the Stability Index Ratio (SIR) comparing EMA-smoothed to raw signal variance per dimension. SIR values uniformly below 1.0 confirm that EMA smoothing reduces volatility across all dimensions, with the greatest stabilization achieved for the clarity dimension (SIR = 0.59).

\begin{table}[htbp]
\centering
\caption{Stability Index Ratio: EMA vs. Raw Signal Variance}
\label{tab:sir}
\begin{tabular}{lcc}
\toprule
\textbf{Dimension} & \textbf{Raw Var.} & \textbf{SIR (EMA/Raw)} \\
\midrule
Valence   & 0.082 & 0.61 \\
Arousal   & 0.074 & 0.63 \\
Stress    & 0.091 & 0.66 \\
Clarity   & 0.068 & 0.59 \\
\midrule
\textbf{Mean} & 0.079 & \textbf{0.62} \\
\bottomrule
\end{tabular}
\end{table}

% ============================================================
\section{Discussion}
\label{sec:discussion}
% ============================================================

The results demonstrate that PSACA achieves stable and coherent psychological state tracking across extended conversational sessions. The high TSS values (mean 0.92) confirm that the adaptive trend modeling pipeline successfully attenuates instantaneous noise while remaining responsive to genuine psychological transitions. The 38.4\% average variance reduction attributed to EMA smoothing highlights the importance of temporal signal processing in affective computing, where raw per-utterance estimates are inherently noisy.

The interaction mode distribution reveals meaningful variation in the system's adaptive behavior. The 18\% incidence of de-escalation mode indicates that the system appropriately identifies psychologically demanding exchanges and modulates its response strategy accordingly. The relatively rare invocation of clarification mode (12\%) suggests that cognitive clarity, while occasionally low, is not the dominant challenge in the evaluated conversational domains.

The selective memory storage policy (MSR = 0.34) reflects a principled trade-off between comprehensive archival and focused retrieval. By storing only the most salient 34\% of turns, the system maintains a memory store that is both informationally dense and computationally tractable. The concentration of storage events around topic transitions and emotionally salient moments (as visualized in Figure~\ref{fig:memory}) validates the psychological salience scoring function of the memory manager.

The emotional shift detection rate of 17\% reflects the rarity of genuine, hysteresis-threshold-exceeding psychological transitions in naturalistic conversation. This rate is consistent with theoretical expectations: most conversational turns represent incremental evolution of psychological state rather than discrete transitions. The detection of shifts at 17\% of turns---concentrated near topic boundaries and mode transitions---confirms that the hysteresis gating mechanism is correctly calibrated to identify substantive rather than spurious transitions.

Trend-conditioned LLM prompting represents a conceptually significant departure from the dominant paradigm of per-turn, state-agnostic prompting. By providing the LLM with a structured representation of the psychological trajectory---including direction of change, recent volatility, and current interaction mode---the system enables response generation that is sensitive to where the user is heading psychologically, not merely where they are at the current turn. This forward-looking conditioning is particularly valuable in de-escalation scenarios, where anticipating continued stress escalation warrants more proactive calming strategies than momentarily high but stable stress.

% ============================================================
\section{Limitations}
\label{sec:limitations}
% ============================================================

Several limitations of the current system warrant acknowledgment.

\textbf{Computational overhead and latency.} The full PSACA pipeline---including ASR, multimodal emotion analysis, fusion, trend computation, memory retrieval, LLM generation, and TTS synthesis---introduces non-trivial computational overhead. The cumulative pipeline latency currently exceeds practical conversational thresholds under resource-constrained deployment conditions. Reducing end-to-end latency through model distillation, asynchronous processing, and hardware acceleration represents a critical direction for future work before deployment in real-time consumer applications.

\textbf{Evaluation scale.} The presented results are based on a single 100-turn evaluation session. While sufficient to characterize system behavior across the principal metrics, this evaluation does not constitute a large-scale user study. Future work should involve multi-session, multi-user longitudinal evaluations with diverse participant populations to assess generalizability, individual differences in affective expression, and long-term memory quality.

\textbf{Ground truth annotation.} The current evaluation relies on algorithmic metrics rather than human annotation of ground-truth psychological states. Establishing gold-standard psychological state labels for conversational turns is methodologically challenging but necessary for rigorous validation of the fusion and trend modeling components. Future work should incorporate expert annotation and self-report psychological state ratings from actual users.

\textbf{Modality fusion calibration.} The confidence-weighted fusion scheme relies on heuristic calibration of modality weights and thresholds. A data-driven approach to weight estimation---for example, through Bayesian fusion or learned attention mechanisms---may yield more principled and generalizable fusion behavior.

\textbf{Cultural and linguistic generalization.} The current system is evaluated on English-language conversational data. Affective expression norms, prosodic conventions, and emotional vocabulary differ substantially across languages and cultures. Extending PSACA to multilingual and cross-cultural settings requires adapted models and evaluation protocols.

% ============================================================
\section{Conclusion}
\label{sec:conclusion}
% ============================================================

We have presented PSACA, a Psychological-State-Aware Conversational AI system that integrates multimodal emotion recognition, adaptive psychological trend modeling, event-driven long-term memory, and trend-conditioned LLM prompting into a cohesive architecture for empathic dialogue. The system tracks a four-dimensional psychological state vector (valence, arousal, stress, clarity) using a layered pipeline of instantaneous fusion, EMA smoothing, regression-based trajectory estimation, and hysteresis gating, achieving Trend Stability Scores between 0.89 and 0.95 across dimensions in a 100-turn evaluation.

Key findings include: (i) EMA-based smoothing reduces psychological state signal variance by 38.4\% on average, substantially improving the quality of state estimates used for LLM conditioning; (ii) the hysteresis-gated shift detection mechanism correctly identifies genuine psychological transitions at a rate of 17\% of turns, concentrated near topic boundaries; (iii) the selective memory storage policy achieves a 34\% storage rate, concentrating archived episodes at psychologically salient moments; and (iv) interaction mode-adaptive TTS synthesis enables the system to modulate its spoken output in response to the inferred psychological context.

PSACA demonstrates that continuous, trajectory-level psychological state tracking is both technically feasible and behaviorally meaningful in real-time conversational AI. By moving beyond per-turn, state-agnostic interaction toward trend-informed, memory-grounded empathic dialogue, this work contributes a step toward conversational AI systems that can engage with the full psychological complexity of human communication. Future work will focus on reducing pipeline latency, expanding evaluation to larger and more diverse user populations, and incorporating human-annotated ground truth for rigorous validation.

% ============================================================
\bibliographystyle{IEEEtran}
% ============================================================

\begin{thebibliography}{99}

\bibitem{picard2000affective}
R. W. Picard, \textit{Affective Computing}. Cambridge, MA, USA: MIT Press, 2000.

\bibitem{schuller2018age}
B. Schuller and A. Batliner, \textit{Computational Paralinguistics: Emotion, Affect and Personality in Speech and Language Processing}. Chichester, UK: Wiley, 2014.

\bibitem{brown2020language}
T. B. Brown et al., ``Language models are few-shot learners,'' in \textit{Proc. Advances in Neural Information Processing Systems (NeurIPS)}, vol. 33, pp. 1877--1901, 2020.

\bibitem{openai2023gpt4}
OpenAI, ``GPT-4 technical report,'' arXiv preprint arXiv:2303.08774, 2023.

\bibitem{morris2018towards}
M. E. Morris, ``Towards an emotional support dialogue system,'' in \textit{Proc. AAAI Workshop on AI for Social Good}, 2018.

\bibitem{fitzpatrick2017delivering}
K. K. Fitzpatrick, A. Darcy, and M. Vierhile, ``Delivering cognitive behavior therapy to young adults with symptoms of depression and anxiety using a fully automated conversational agent (Woebot): A randomized controlled trial,'' \textit{JMIR Mental Health}, vol. 4, no. 2, p. e19, 2017.

\bibitem{poria2017review}
S. Poria, E. Cambria, R. Bajpai, and A. Hussain, ``A review of affective computing: From unimodal analysis to multimodal fusion,'' \textit{Information Fusion}, vol. 37, pp. 98--125, 2017.

\bibitem{sebe2005multimodal}
N. Sebe, I. Cohen, and T. Gevers, ``Multimodal approaches for emotion recognition: A survey,'' \textit{Proc. SPIE}, vol. 5670, pp. 56--67, 2005.

\bibitem{cowie2001emotion}
R. Cowie et al., ``Emotion recognition in human-computer interaction,'' \textit{IEEE Signal Processing Magazine}, vol. 18, no. 1, pp. 32--80, 2001.

\bibitem{zhong2022dialoglm}
P. Zhong et al., ``DialogLM: Pre-trained model for long dialogue understanding and summarization,'' in \textit{Proc. AAAI Conference on Artificial Intelligence}, vol. 36, pp. 11765--11773, 2022.

\bibitem{ekman1978facial}
P. Ekman and W. Friesen, \textit{Facial Action Coding System}. Palo Alto, CA, USA: Consulting Psychologists Press, 1978.

\bibitem{schuller2003hidden}
B. Schuller, G. Rigoll, and M. Lang, ``Hidden Markov model-based speech emotion recognition,'' in \textit{Proc. IEEE International Conference on Acoustics, Speech, and Signal Processing (ICASSP)}, vol. 2, pp. II-1--II-4, 2003.

\bibitem{healey2005detecting}
J. A. Healey and R. W. Picard, ``Detecting stress during real-world driving tasks using physiological sensors,'' \textit{IEEE Transactions on Intelligent Transportation Systems}, vol. 6, no. 2, pp. 156--166, 2005.

\bibitem{zadeh2018multimodal}
A. Zadeh, P. P. Liang, N. Mazumder, S. Poria, E. Cambria, and L.-P. Morency, ``Memory fusion network for multi-view sequential learning,'' in \textit{Proc. AAAI Conference on Artificial Intelligence}, vol. 32, 2018.

\bibitem{tsai2019multimodal}
Y.-H. H. Tsai, S. Bai, P. P. Liang, J. Z. Kolter, L.-P. Morency, and R. Salakhutdinov, ``Multimodal transformer for unaligned multimodal language sequences,'' in \textit{Proc. 57th Annual Meeting of the Association for Computational Linguistics (ACL)}, pp. 6558--6569, 2019.

\bibitem{devlin2019bert}
J. Devlin, M.-W. Chang, K. Lee, and K. Toutanova, ``BERT: Pre-training of deep bidirectional transformers for language understanding,'' in \textit{Proc. NAACL-HLT}, pp. 4171--4186, 2019.

\bibitem{sun2019utilizing}
C. Sun, X. Qiu, Y. Xu, and X. Huang, ``How to fine-tune BERT for text classification?'' in \textit{Proc. CCF International Conference on Natural Language Processing and Chinese Computing}, pp. 194--206, 2019.

\bibitem{eyben2016geneva}
F. Eyben et al., ``The Geneva minimalistic acoustic parameter set (GeMAPS) for voice research and affective computing,'' \textit{IEEE Transactions on Affective Computing}, vol. 7, no. 2, pp. 190--202, 2016.

\bibitem{russell1980circumplex}
J. A. Russell, ``A circumplex model of affect,'' \textit{Journal of Personality and Social Psychology}, vol. 39, no. 6, pp. 1161--1178, 1980.

\bibitem{mehrabian1996pleasure}
A. Mehrabian, ``Pleasure-arousal-dominance: A general framework for describing and measuring individual differences in temperament,'' \textit{Current Psychology}, vol. 14, no. 4, pp. 261--292, 1996.

\bibitem{sweller1988cognitive}
J. Sweller, ``Cognitive load during problem solving: Effects on learning,'' \textit{Cognitive Science}, vol. 12, no. 2, pp. 257--285, 1988.

\bibitem{rabiner1993fundamentals}
L. Rabiner and B.-H. Juang, \textit{Fundamentals of Speech Recognition}. Englewood Cliffs, NJ, USA: Prentice Hall, 1993.

\bibitem{metallinou2013tracking}
A. Metallinou and S. Narayanan, ``Tracking changes in continuous emotion states using body language and prosodic cues,'' in \textit{Proc. IEEE International Conference on Acoustics, Speech and Signal Processing (ICASSP)}, pp. 3581--3585, 2013.

\bibitem{sukhbaatar2015end}
S. Sukhbaatar, A. Szlam, J. Weston, and R. Fergus, ``End-to-end memory networks,'' in \textit{Proc. Advances in Neural Information Processing Systems (NeurIPS)}, vol. 28, 2015.

\bibitem{lewis2020retrieval}
P. Lewis et al., ``Retrieval-augmented generation for knowledge-intensive NLP tasks,'' in \textit{Proc. Advances in Neural Information Processing Systems (NeurIPS)}, vol. 33, pp. 9459--9474, 2020.

\bibitem{karpukhin2020dense}
V. Karpukhin et al., ``Dense passage retrieval for open-domain question answering,'' in \textit{Proc. 2020 Conference on Empirical Methods in Natural Language Processing (EMNLP)}, pp. 6769--6781, 2020.

\bibitem{xu2021beyond}
C. Xu et al., ``Beyond goldfish memory: Long-term open-domain conversation,'' arXiv preprint arXiv:2107.07567, 2021.

\bibitem{rashkin2019towards}
H. Rashkin, E. M. Smith, M. Li, and Y.-L. Boureau, ``Towards empathetic open-domain conversation models: A new benchmark and dataset,'' in \textit{Proc. 57th Annual Meeting of the Association for Computational Linguistics (ACL)}, pp. 5370--5381, 2019.

\bibitem{majumder2020mime}
N. Majumder et al., ``MIME: MIMicking emotions for empathetic response generation,'' in \textit{Proc. 2020 Conference on Empirical Methods in Natural Language Processing (EMNLP)}, pp. 8968--8979, 2020.

\bibitem{shen2018natural}
J. Shen et al., ``Natural TTS synthesis by conditioning WaveNet on mel spectrogram predictions,'' in \textit{Proc. IEEE International Conference on Acoustics, Speech, and Signal Processing (ICASSP)}, pp. 4779--4783, 2018.

\bibitem{ren2020fastspeech}
Y. Ren et al., ``FastSpeech 2: Fast and high-quality end-to-end text to speech,'' in \textit{Proc. International Conference on Learning Representations (ICLR)}, 2021.

\bibitem{lei2022msemotts}
Y. Lei, S. Yang, X. Wang, and L. Xie, ``MSEmoTTS: Multi-scale emotion transfer, prediction, and control for emotional speech synthesis,'' \textit{IEEE/ACM Transactions on Audio, Speech, and Language Processing}, vol. 30, pp. 853--864, 2022.

\end{thebibliography}

\end{document}